\section{Discussion, Limitations, and Conclusion}
\label{sec:conclusion}

\paragraph{Discussion.} Our results argue that the two prevailing programs for
deployment-time improvement---accumulating experience outside the model and
optimizing the scalar feedback bit with policy gradients---leave value on the
table that a third route captures: internalizing each verified experience into
the weights through anchored self-distillation. The streaming schedule is not an
obstacle to engineer around but an asset: it supplies an emergent curriculum and
a data flywheel that batch training on the same data cannot reproduce.

\paragraph{Limitations.} Our experiments use a single 7B code-centric model with
LoRA adapters and one stream order per task (StreamBench's released sequences);
multi-seed variance is not characterized. Feedback is binary---richer textual
feedback \citep{hue2026sdpo} is unexplored in our setting. Our method spends
${\sim}10\times$ more oracle calls than the baselines (on memory maintenance,
not scoring), a real cost where verification is expensive. General-capability
retention beyond the streamed tasks (e.g., on knowledge benchmarks) is not yet
measured. \todo{Transfer-to-unseen-domain runs and the exploration-matched batch
arm are pending.}

\paragraph{Conclusion.} We presented \method{}, to our knowledge the first
method to perform verified, weight-updating self-distillation online over a
deployment stream: verified self-generated answers feed a reward-filtered
memory that serves retrieval today and distillation forever, with forward
re-evaluation converting weight gains back into training data. \method{} is
best on all seven StreamBench tasks, absorbs a three-domain distribution shift
that collapses a stability-engineered RL baseline---with one model beating
per-domain specialists through weight-borne transfer---and outperforms batch
training with the identical objective on identical data. Deployed models need
not be frozen: a single pass of grounded feedback is enough to learn from.

\subsection*{Reproducibility Statement}
All benchmarks are public (StreamBench and its constituent datasets;
Table~\ref{tab:datasets} lists sources), and all methods use an open-weights
model (Qwen2.5-Coder-7B-Instruct) with fixed released stream orders (seed 42),
making every run exactly repeatable. Complete hyperparameters appear in
Appendix~\ref{app:hparams}, prompt templates in Appendix~\ref{app:prompts}, and
oracle definitions in \S\ref{sec:exp-setup}. Our released repository contains
all training and evaluation code, per-problem logs for every run in every
table, the table-generation scripts that produce the paper's numbers directly
from those logs, and the causal-ordering audit of Appendix~\ref{app:leakage}.
